\chapter{项目三：配置加载器}

\section{项目目标}

这个项目实现一个小型配置加载器。输入是 \code{key=value} 文件，支持注释、空行、重复 key 覆盖、类型转换和必填项校验。输出是一个强类型配置对象。

示例配置：

\begin{lstlisting}[caption={配置文件示例}]
# server config
host = 127.0.0.1
port = 8080
workers = 4
log_level = info
\end{lstlisting}

目标不是写一个通用配置库，而是把字符串处理、错误接口、类型建模、文件 I/O 和测试组织合在一起。配置加载器是非常适合练基本功的项目，因为它看起来简单，但边界很多。

\section{先定义业务配置}

不要让整个程序到处传 \code{unordered_map<string, string>}。读取阶段可以用字符串映射，业务阶段应该使用强类型：

\begin{lstlisting}[language=C++,caption={强类型服务器配置}]
enum class LogLevel {
    debug,
    info,
    warning,
    error,
};

struct ServerConfig {
    std::string host;
    int port = 0;
    int workers = 0;
    LogLevel log_level = LogLevel::info;
};
\end{lstlisting}

这一步把业务要求显式化：\code{port} 是整数，\code{workers} 是整数，\code{log_level} 只能是枚举值。后续代码不需要反复解析字符串。

\section{错误类型要能定位问题}

配置错误需要带上下文：

\begin{lstlisting}[language=C++,caption={配置错误}]
enum class ConfigErrorKind {
    file_not_found,
    invalid_line,
    missing_key,
    invalid_value,
};

struct ConfigError {
    ConfigErrorKind kind = ConfigErrorKind::invalid_line;
    int line = 0;
    std::string key;
    std::string message;
};

struct ConfigResult {
    std::optional<ServerConfig> config;
    std::optional<ConfigError> error;
};
\end{lstlisting}

为什么要有 \code{line} 和 \code{key}？因为用户修配置时需要知道哪里错。错误结构不是为了程序员自娱自乐，而是为了让调用者能做出具体动作。

\section{解析阶段：从行到映射}

先把文件解析成键值映射：

\begin{lstlisting}[language=C++,caption={解析配置条目}]
using RawConfig = std::unordered_map<std::string, std::string>;

struct RawConfigResult {
    RawConfig values;
    std::optional<ConfigError> error;
};

RawConfigResult parse_raw_config(std::istream& input)
{
    RawConfigResult result;
    std::string line;
    int line_number = 0;
    while (std::getline(input, line)) {
        ++line_number;
        try {
            std::optional<ConfigEntry> entry = parse_config_line(line);
            if (entry.has_value()) {
                result.values[entry->key] = entry->value;
            }
        } catch (const std::exception& error) {
            return {.error = ConfigError{
                .kind = ConfigErrorKind::invalid_line,
                .line = line_number,
                .message = error.what(),
            }};
        }
    }
    return result;
}
\end{lstlisting}

重复 key 的规则在这里定义：后者覆盖前者。真实项目也可以记录警告，但规则必须写清楚，不能让行为由容器偶然决定。

\section{类型转换阶段}

字符串转整数要检查完整消费和范围：

\begin{lstlisting}[language=C++,caption={解析整数配置}]
std::optional<int> parse_int_value(std::string_view text)
{
    int value = 0;
    const char* begin = text.data();
    const char* end = text.data() + text.size();
    const auto [ptr, ec] = std::from_chars(begin, end, value);
    if (ec != std::errc{} || ptr != end) {
        return std::nullopt;
    }
    return value;
}
\end{lstlisting}

日志级别是有限集合：

\begin{lstlisting}[language=C++,caption={解析日志级别}]
std::optional<LogLevel> parse_log_level(std::string_view text)
{
    if (text == "debug") {
        return LogLevel::debug;
    }
    if (text == "info") {
        return LogLevel::info;
    }
    if (text == "warning") {
        return LogLevel::warning;
    }
    if (text == "error") {
        return LogLevel::error;
    }
    return std::nullopt;
}
\end{lstlisting}

用 \code{enum class} 的好处是后续 \code{switch} 不会把日志级别和普通整数混用。

\section{从映射构造强类型配置}

构造 \code{ServerConfig} 时处理必填项、默认值和范围：

\begin{lstlisting}[language=C++,caption={构造强类型配置}]
std::optional<std::string_view> find_raw(const RawConfig& raw, std::string_view key)
{
    const auto found = raw.find(std::string{key});
    if (found == raw.end()) {
        return std::nullopt;
    }
    return std::string_view{found->second};
}

ConfigResult build_server_config(const RawConfig& raw)
{
    ServerConfig config;

    const std::optional<std::string_view> host = find_raw(raw, "host");
    if (!host.has_value()) {
        return {.error = ConfigError{.kind = ConfigErrorKind::missing_key,
                                     .key = "host",
                                     .message = "missing host"}};
    }
    config.host = *host;

    const std::optional<std::string_view> port_text = find_raw(raw, "port");
    const std::optional<int> port = port_text.has_value()
        ? parse_int_value(*port_text)
        : std::nullopt;
    if (!port.has_value() || *port < 1 || *port > 65535) {
        return {.error = ConfigError{.kind = ConfigErrorKind::invalid_value,
                                     .key = "port",
                                     .message = "port must be 1..65535"}};
    }
    config.port = *port;

    config.workers = 1;
    if (const std::optional<std::string_view> workers = find_raw(raw, "workers")) {
        const std::optional<int> parsed = parse_int_value(*workers);
        if (!parsed.has_value() || *parsed <= 0) {
            return {.error = ConfigError{.kind = ConfigErrorKind::invalid_value,
                                         .key = "workers",
                                         .message = "workers must be positive"}};
        }
        config.workers = *parsed;
    }

    if (const std::optional<std::string_view> level = find_raw(raw, "log_level")) {
        const std::optional<LogLevel> parsed = parse_log_level(*level);
        if (!parsed.has_value()) {
            return {.error = ConfigError{.kind = ConfigErrorKind::invalid_value,
                                         .key = "log_level",
                                         .message = "invalid log level"}};
        }
        config.log_level = *parsed;
    }

    return {.config = std::move(config)};
}
\end{lstlisting}

这段代码刻意分阶段：先把文本读成原始映射，再把映射转换成强类型对象。这样解析行格式和业务校验可以分别测试。

\section{文件入口}

最终文件入口很薄：

\begin{lstlisting}[language=C++,caption={加载配置文件}]
ConfigResult load_config_file(const std::filesystem::path& path)
{
    std::ifstream input{path};
    if (!input) {
        return {.error = ConfigError{.kind = ConfigErrorKind::file_not_found,
                                     .message = "cannot open config file"}};
    }

    RawConfigResult raw = parse_raw_config(input);
    if (raw.error.has_value()) {
        return {.error = raw.error};
    }
    return build_server_config(raw.values);
}
\end{lstlisting}

这里没有把所有逻辑塞进一个函数。文件打开、行解析、类型转换各有边界，错误也能保留上下文。

\section{测试计划}

这个项目至少要覆盖：

\begin{itemize}
  \item 空行和注释行被忽略。
  \item 缺少等号返回行号错误。
  \item 重复 key 使用后者。
  \item 缺少必填 \code{host} 或 \code{port} 报错。
  \item \code{port=abc}、\code{port=0}、\code{port=65536} 报错。
  \item \code{workers} 缺失时使用默认值 \code{1}。
  \item 非法 \code{log_level} 报错。
\end{itemize}

测试尽量用 \code{std::istringstream}，不要每个单元测试都创建真实文件。只有 \code{load_config_file} 这一层需要文件系统测试。

\begin{exercisebox}
给配置加载器增加 \code{include=other.conf}。要求避免递归实现，使用显式队列或栈处理待加载文件，并检测重复包含。写清楚相对路径以哪个目录为基准。
\end{exercisebox}
